% =============================================================================
%  preamble.tex -- WCA-Bench 中文版宏包与宏定义
%
%  设计约束
%  --------
%   * 必须使用 XeLaTeX 编译（ctexart 默认调用 xeCJK）。
%   * 仅使用 CTAN 常见宏包，标准 TeX Live 2024 可直接编译。
%   * 中文排版由 ctex 处理（默认 fontset=auto，在 Windows 上自动使用系统
%     中文字体，如 SimSun；在其他平台自动选择可用中文字体）。
%   * 专有名词（WCA-Bench、Plackett--Luce、XGBoost、GNN、DNF、
%     Psych Sheet 等）保留英文。
% =============================================================================

% --- 版面 -------------------------------------------------------------------
\usepackage[margin=1in]{geometry}
\setlength{\emergencystretch}{3em}

% --- 数学 -------------------------------------------------------------------
\usepackage{amsmath}
\usepackage{amssymb}
\usepackage{mathtools}

% --- 表格 -------------------------------------------------------------------
\usepackage{booktabs}
\usepackage{multirow}
\usepackage{array}
\usepackage{tabularx}
\usepackage{threeparttable}
\usepackage{makecell}

% --- 图 ---------------------------------------------------------------------
\usepackage{graphicx}
\usepackage{xcolor}
\usepackage{caption}
\usepackage{subcaption}
\graphicspath{{../figures/}}

% --- 算法 -------------------------------------------------------------------
\usepackage{algorithm}
\usepackage{algpseudocode}
\floatname{algorithm}{算法}

% --- 排版与列表 -------------------------------------------------------------
\usepackage{microtype}
\usepackage{enumitem}
\setlist[itemize]{leftmargin=1.6em,topsep=2pt,itemsep=1pt}
\setlist[enumerate]{leftmargin=1.8em,topsep=2pt,itemsep=1pt}

% --- 超链接与文献（置于最后） ----------------------------------------------
\usepackage{url}
\usepackage[numbers,sort&compress]{natbib}
\usepackage[unicode,hidelinks,breaklinks]{hyperref}

% --- 图注 -------------------------------------------------------------------
\captionsetup{font=small,labelfont=bf}

% --- 颜色 -------------------------------------------------------------------
\definecolor{wcablu}{HTML}{1F4E79}
\definecolor{wcared}{HTML}{B03A2E}
\definecolor{wcagrn}{HTML}{1E6B52}

% =============================================================================
%  记号宏
% =============================================================================
\newcommand{\wca}{WCA}
\newcommand{\wcabench}{WCA-Bench}
\newcommand{\task}[1]{\textsf{T#1}}
\newcommand{\x}{\mathbf{x}}
\newcommand{\y}{y}
\newcommand{\asof}{\textsf{as\_of}}
\newcommand{\metric}{\mathcal{M}}
\newcommand{\loss}{\mathcal{L}}
\newcommand{\DNF}{\textsc{DNF}}
\newcommand{\DNS}{\textsc{DNS}}
\newcommand{\kendall}{\ensuremath{\tau}}
\newcommand{\brier}{\ensuremath{B}}
\newcommand{\aucpr}{\ensuremath{\mathrm{AUC\text{-}PR}}}
\newcommand{\aucroc}{\ensuremath{\mathrm{AUC\text{-}ROC}}}
\newcommand{\maelog}{\ensuremath{\mathrm{MAE}_{\log}}}
\newcommand{\rmselog}{\ensuremath{\mathrm{RMSE}_{\log}}}
\newcommand{\ev}{\text{event}}
\newcommand{\code}[1]{\texttt{#1}}

% 表格中“越小越好 / 越大越好”的箭头
\newcommand{\arrdown}{$\downarrow$}
\newcommand{\arrup}{$\uparrow$}

% =============================================================================
%  算法环境中文关键字
% =============================================================================
\algrenewcommand\algorithmicrequire{\textbf{输入：}}
\algrenewcommand\algorithmicensure{\textbf{输出：}}
\algrenewcommand\algorithmicif{如果}
\algrenewcommand\algorithmicthen{则}
\algrenewcommand\algorithmicelse{否则}
\algrenewcommand\algorithmicfor{对于}
\algrenewcommand\algorithmicdo{执行}
\algrenewcommand\algorithmicreturn{返回}
\algrenewcommand\algorithmicwhile{当}
\algrenewcommand\algorithmicend{结束}

\endinput
